\documentclass[../candidatura.tex]{subfiles}
\graphicspath{{\subfix{../imagenes/}}}
\begin{document}

\section{Metodología}

% en borradores iniciales, la extensión sugerida era de 2 páginas

Esta tesis es una investigación desde la práctica artística, entendida como un lugar de producción de conocimiento desde lo visceral, lo corporal, desde el cuerpo en escena, desde el diálogo y retroalimentación entre artista y sus materiales de trabajo. Desde el aserrín de este taller constante emergen prototipos de software y hardware que permiten la síntesis de sintetizadores, cada prototipo siendo a la vez una obra de arte, una bitácora de experimentos, un rastro de decisiones de diseño, y una estrella dentro de una constelación de herramientas que permiten nuevas formas de hacer ruido.

La relación entre artista, materiales y obra (en este caso fabricante, circuito y sintetizador) ha sido descrita por el diseñador de instrumentos Don Buchla en un poema de 1971, donde propone que las máquinas sonoras pueden permitir y producir un nuevo lenguaje y nuevos sentidos de escucha \footcite[88--89]{gordon_composers_2025}.

Se propone también la metáfora de los sintetizadores como patinetas: esta tesis trata sobre las inevitables fuerzas gravitacionales, sobre materiales y estrategias de construcción de patinetas como rodamientos, ruedas y maderas, y por sobre todo es la construcción de un skatepark donde poder proponer nuevas piruetas, o más bien, ruidos y modulaciones.

% revisar JJ Gibson, affordances, que aparece en el poema original

\subsection{Análisis de referentes}

Esta tesis incluye una revisión bibliográfica sobre referentes artísticos, académicos, comerciales, de hardware y software abierto y cerrado, de electrónica y de computación aplicada a artes y diseño.

Esta revisión se enfocará en temáticas afines a la síntesis de sintetizadores de forma popular, para así identificar conceptos y debates relevantes que permitan contextualizar la investigación y práctica artística realizada durante esta investigación doctoral.

\subsection{Prototipado electrónico}

La práctica electrónica se hará de forma iterativa y rumiante, creando progresivamente distintos prototipos de módulos constituyentes de sintetizadores. Su socialización y su uso permitirán el aprendizaje a partir de errores, la escritura de estándares y estrategias de fabricación y uso, que son una práctica común en la comunidad de hardware libre, y que permitirá la creación de sintetizadores con distintas variantes.

Se seguirán preceptos industriales y comerciales, de diseño a partir de componentes discretos, de montaje en placas de prototipado, de modelamiento computacional, de fabricación de placas como soporte electromecánico y de soldadura, y de lanzamiento de distintas revisiones, pasando de tecnología THT a montaje de superficie para abaratar costos en horas humano de soldadura.

\subsection{Prototipado computacional}

La práctica computacional se hará de forma iterativa y rumiante, publicando productos de viabilidad mínima, y compartiendo el proceso de desarrollo en repositorios públicos, con el fin de que puedan todos los artefactos de la investigación ser analizados, replicados, adaptados y mejorados por pares artistas, académicas, y audiencias interesadas, a través del tiempo y de forma descentralizada.

Se realizarán primero productos de viabilidad mínima en distintos lenguajes, y una vez se logre destilar lo esencial de lo necesario en el código, se reescribirá el código en lenguajes de menor nivel como C++ y Rust para así optimizarlo y permitir su uso en dispositivos de bajo consumo como microcontroladores. También se estima importante una implementación modelada en código usando la herramienta gratuita VCV Rack \footcite{gabrielli_developing_2020}.

Este trabajo está sintetizado en la colección de Popusintes \footcite{piruetas_popusintesrack_2026}, y en las distintas bibliotecas de elementos interfásicos para su uso con microcontroladores, para botones \footcite{piruetas_boton_2026}, perillas \footcite{piruetas_perilla_2026} y micrófonos \footcite{piruetas_microfono_2026}, entre otros.

\subsection{Escritura de documentación}

La documentación se hará de forma paralela al desarrollo de electrónica y computación, y permitirá que el proyecto de popusintesíntesis sea legible y pueda ser testeado y validado por agentes externos. No basta con que los módulos y los sintetizadores funcionen, también debe haber claridad en sus decisiones de diseño, su contexto de uso, su implementación, y cómo se pueden modificar sus parámetros para otros fines.

\subsection{Socialización de la investigación}

Durante la investigación doctoral se desarrollarán distintas herramientas que faciliten la síntesis de sintetizadores, las que serán prototipadas, documentadas y puestas a disposición de pares artistas y académicas, y serán usadas en contextos de producción de obra artística, como material docente para la enseñanza de electrónica y computación, y posteriormente serán comercializadas para difundir y financiar la investigación material realizada.

Se considera relevante para el desarrollo correcto de esta investigación la validación de otras artistas que pongan en práctica el uso de las herramientas lanzadas durante el desarrollo de la tesis.

\subsection{Escritura y enseñanza de material base}

Para poder abordar de forma sostenida una investigación doctoral profunda en la síntesis de sintetizadores, se están desarrollando en paralelo cursos de pregrado en diseño que permitan a estudiantes poner en práctica lo investigado, y que esto a su vez nutra de nuevas ideas y experiencias a esta investigación. Como subproducto de esta enseñanza se están escribiendo libros de referencia de enseñanza de electrónica y computación aplicadas a sonido, que permitan un diálogo constante con audiencias de estudiantes, artistas e investigadoras.

\end{document}

